
\section{\BookRecordIndex{T0054}伴随不等式}\label{sec:2-6}
\zhindex{bansuibudengshi}{伴随不等式}

本章最后讨论殆复四维流形中 \BookRecordIndex{T0003}$J$-全纯曲线的\BookRecordIndex{T0054}伴随不等式。证明将在附录 E 给出。作为热身，
下面的习题考察 \BookRecordIndex{T0003}$J$-全纯曲线与实余维为二的复子流形之间的交点。

\begin{exercise}\label{ex:2-6-1}
设 $Q$ 是 $(M,J)$ 的紧致实余维二子流形，并且在 $JTQ=TQ$ 的意义下是 $J$-全纯的；
又设 $u:B_1\to(M,J)$ 是一条满足 $u(0)\in Q$ 的 \BookRecordIndex{T0003}$J$-全纯曲线。
\begin{enumerate}[label=(\roman*),leftmargin=3em]
  \item 证明这个交点是孤立的，除非 $u(B_1)\subset Q$。

  \textit{提示：}仿照引理~\ref{lem:2-4-2} 找到合适的局部坐标，再模仿
  引理~\ref{lem:2-4-3} 的证明。

  \item 假设 $u(B_1)\not\subset Q$。必要时缩小 $B_1$，可以假设 $u(0)$ 是
  $u(B_1)$ 与 $Q$ 相交的唯一一点。把 $u$ 与 $Q$ 的\BookRecordIndex{T0201}局部交数 $u\cdot Q$ 定义为
  $v$ 与 $Q$ 的交点个数（计入重数），其中
  $v:B_1\to M$ 是一个光滑映射\sourcefootnote{原书把 $v$ 的靶空间印作 $Q$；然而这里须计算 $v$ 与
  $Q$ 的交点，且 $v$ 是保持边界值的 $u:B_1\to M$ 的扰动，故改为 $M$。}；
  它是 $u$ 在边界固定的光滑映射空间中的一个通有扰动，即
  $v|_{\partial B_1}=u|_{\partial B_1}$。证明 $u\cdot Q\ge1$，
  且等号成立当且仅当 $u$ 在零点与 $Q$ 横截。
\end{enumerate}
\end{exercise}

上述习题是\emph{\BookRecordIndex{T0186}交点正性}最简单的情形；这一现象最早由
\personlater{geluomofu}{格罗莫夫}{Gromov}在 \srcite{160} 中注意到。
它在四维情形尤其重要，因为通过观察一条曲线与其他曲线以及与自身的交点，可以获得有关
曲线奇点的信息。习题~\ref{ex:2-6-1} 表明，在四维中，两条 \BookRecordIndex{T0003}$J$-全纯曲线的一个交点，
只要在其中一条曲线上是\BookRecordIndex{T0121}非奇异点，就对交数作正贡献。对这一论证作一个相当不平凡的细化，
可知每个交点都作正贡献，而且曲线 $C$ 上的每个奇点都对（局部）自交数 $C\cdot C$
作正贡献。后一结果导出下面定理~\ref{thm:2-6-4} 中的\BookRecordIndex{T0055}伴随公式；前一结果陈述为
定理~\ref{thm:2-6-3}。两个定理用相同的方法证明，但二者互不蕴含。证明将在附录 E 给出。

\begin{definition}\label{def:2-6-2}
设 $(M,J)$ 是殆复四维流形；对 $i=0,1$，设 $(\Sigma_i,j_i)$ 是闭黎曼曲面，并设
$u_i:\Sigma_i\to M$ 是 \BookRecordIndex{T0003}$J$-全纯曲线。$u_0$ 与 $u_1$ 的一个\emph{\BookRecordIndex{T0165}横截交点}，
是满足下式的对子 $(z_0,z_1)\in\Sigma_0\times\Sigma_1$：
\begin{equation}\label{eq:2-6-1}
  u_0(z_0)=u_1(z_1)=:x,
  \qquad
  T_xM=\operatorname{im}du_0(z_0)\oplus\operatorname{im}du_1(z_1).
\end{equation}
在 $\Sigma_0=\Sigma_1=\Sigma$ 且 $u_0=u_1=u$ 的情形，这样的对子 $(z_0,z_1)$
称为 $u$ 的一个\emph{\BookRecordIndex{T0167}横截自交点}。
\end{definition}

所有交点的数目记为
\[
  \delta(u_0,u_1)
  :=\#\bigl\{(z_0,z_1)\in\Sigma_0\times\Sigma_1
  \bigm|u_0(z_0)=u_1(z_1)\bigr\},
\]
而曲线 $u$ 的所有自交点的数目记为
\[
  \delta(u)
  :=\frac12\#\bigl\{(z_0,z_1)\in\Sigma\times\Sigma
  \bigm|z_0\ne z_1,\ u(z_0)=u(z_1)\bigr\}.
\]
显然，只要 $u$ 不是单纯的，这个数就是无穷；同样，当两条曲线的并集不是单纯的时，
$\delta(u_0,u_1)$ 也是无穷。在单纯情形，这些数为有限这一点并不能直接从定义看出。
在这种情形下，引理~\ref{lem:2-4-3} 表明，$\Sigma_0\times\Sigma_1$ 中的交点只能在
由两条曲线的奇点组成的点对处聚积。不过，下面的定理~\ref{thm:2-6-3} 和
定理~\ref{thm:2-6-4}
给出了单纯情形下 $\delta(u_0,u_1)$ 与 $\delta(u)$ 的拓扑上界。记
$c_1(A)=\langle c_1(TM,J),A\rangle$ 为二维同调类 $A$ 的第一陈数，
以 $A_0\cdot A_1$ 表示两个类 $A_0,A_1$ 的交数，
以 $\chi(\Sigma)$ 表示闭黎曼曲面 $\Sigma$ 的欧拉示性数。下面两个定理中的黎曼曲面都是
闭的，但未必连通。

\begin{theorem}[\BookRecordIndex{T0186}交点正性]\label{thm:2-6-3}
设 $(M,J)$ 是殆复四维流形，$A_0,A_1\in H_2(M;\mathbb Z)$ 是分别由单纯
\BookRecordIndex{T0003}$J$-全纯曲线 $u_0:\Sigma_0\to M$ 与 $u_1:\Sigma_1\to M$ 表示的同调类。
假设两条曲线的并集是单纯的（即对任意两个非空开子集
$U_0\subset\Sigma_0$ 与 $U_1\subset\Sigma_1$，都有
$u_0(U_0)\ne u_1(U_1)$）。则
\[
  \delta(u_0,u_1)\le A_0\cdot A_1,
\]
而且等号成立当且仅当所有交点都横截。
\end{theorem}

\begin{proof}
见第 E.1 节和第 E.2 节。
\end{proof}

\begin{theorem}[\BookRecordIndex{T0054}伴随不等式]\label{thm:2-6-4}
设 $(M,J)$ 是殆复四维流形，$A\in H_2(M;\mathbb Z)$ 是由一条单纯 \BookRecordIndex{T0003}$J$-全纯曲线
$u:\Sigma\to M$ 表示的同调类。则
\begin{equation}\label{eq:2-6-2}
  2\delta(u)-\chi(\Sigma)\le A\cdot A-c_1(A).
\end{equation}
等号成立当且仅当 $u$ 是仅有\BookRecordIndex{T0167}横截自交点的浸入。
\end{theorem}

\begin{proof}
见第 E.1 节和第 E.2 节。
\end{proof}

一个直接推论是：嵌入 \BookRecordIndex{T0003}$J$-全纯曲线序列 $u_\nu:\Sigma\to M^4$ 不可能收敛到一条
带有奇点或自交点的\BookRecordIndex{T0091}单纯曲线 $u:\Sigma\to M^4$。

\begin{example}\label{ex:2-6-5}
设 $J_0$ 是 $\mathbb{CP}^2$ 上的标准复结构，$\Sigma$ 是亏格为 $g$ 的连通黎曼曲面。
则一条次数为 $d$ 的全纯曲线 $u:\Sigma\to\mathbb{CP}^2$ 表示某个类 $A$，且
$A\cdot A=d^2$、$c_1(A)=3d$。因此\BookRecordIndex{T0054}伴随不等式具有形式
\[
  \delta+g\le\frac{(d-1)(d-2)}2.
\]
一个不连通的例子是 $\mathbb{CP}^2$ 中 $d$ 条互异直线的并。此时 $\Sigma$ 有 $d$ 个
连通分支，每个分支的亏格都是零，因而 $\chi(\Sigma)=2d$。由于
$2\delta=d(d-1)$，不等式两边相等。若其中两条直线重合，则原书给出的数值为
$2\delta=d(d-1)-2$，\sourcefootnote{原书此处把两条分量重合时的 $2\delta$ 写成
有限值；但按本节在定义~\ref{def:2-6-2} 后给出的 $\delta(u)$ 定义及紧接的说明，
非\BookRecordIndex{T0091}单纯曲线的 $\delta(u)$ 为无穷，与此不\BookRecordIndex{T0348}相容。译文照录原式，不自行重构作者意图。}
所以不等式
\eqref{eq:2-6-2}
仍然成立（且为严格不等式），
尽管 $u$ 并不单纯。
\end{example}

\begin{exercise}\label{ex:2-6-6}
设 $(M,J)$ 是殆复四维流形，并且 $M$ 中每条 \BookRecordIndex{T0003}$J$-全纯曲线都有正自交数。证明不等式
$c_1(A)-A\cdot A\le\chi(\Sigma)$ 对由非\BookRecordIndex{T0091}单纯曲线表示的类仍然成立；而且除非每条
表示类 $A$ 的 \BookRecordIndex{T0003}$J$-全纯曲线
$u:\Sigma\to M$ 都是嵌入，否则该不等式是严格的。找出一条满足 $A\cdot A<0$ 且违反
这一不等式的非\BookRecordIndex{T0091}单纯曲线。
\end{exercise}

\begin{exercise}\label{ex:2-6-7}
证明：当 $u_0$ 与 $u_1$ 是相互横截的浸入，并且仅有\BookRecordIndex{T0167}横截自交点时，定理~\ref{thm:2-6-3}
的断言可由定理~\ref{thm:2-6-4} \BookRecordIndex{T0326}推出。
\end{exercise}

\BookRecordIndex{T0003}$J$-全纯曲线的\BookRecordIndex{T0054}伴随不等式最早由 McDuff 在 \srcite{256} 中陈述。Micallef 指出，其中一个
引理错误地选取了坐标，从而影响了后续许多细节。McDuff 在 \srcite{264} 中给出了修正并
简化的版本。该论证表明，若 $z$ 是 $u$ 的\BookRecordIndex{T0233}临界点，则 $u$ 的任何局部
$J$-全纯扰动，只要它是浸入，就必定至少有一个\BookRecordIndex{T0116}二重点。另一种办法是分析像曲线
$u(\Sigma)$ 的奇点拓扑类型
（这里无需把曲线限制在四维空间中）。McDuff 曾在 \srcite{261} 中尝试这一方法，
但遗憾的是，该论文含有一些严重错误，并且命题 2.8 的证明并不成立。即使把它修正，
这种方法基本上也是拓扑性的，至多只能在 $C^0$-等价意义下对这些芽作分类。
在 \srcite{287} 中，Micallef 和 White 成功证明了 \BookRecordIndex{T0003}$J$-全纯曲线上任意奇点的芽都在
$C^1$ 意义下等价于一个标准全纯芽；
这是一个关于\BookRecordIndex{T0175}极小曲面奇点的远为一般的结果的特殊情形。Sikorav 在 \srcite{379} 中对这些
问题给出了更直接的处理，并对 $J$ 只作尽可能弱的光滑性假设。在附录 E 中，Lazzarini
描述了 Micallef--White 论证的另一版本，它对光滑 \BookRecordIndex{T0003}$J$-全纯曲线成立。

一个引起关注的相关问题是：$\mathbb{CP}^2$ 中任意辛嵌入的闭二维流形是否都\BookRecordIndex{T0361}辛同痕于一条
复曲线。关于这个问题的一些进展及参考文献目录，见 Shevchishin \srcite{375}。
